% Source PDF: source/普通高考/2013/2013安徽理.pdf, page 1.
% PDF embedded bitmap attempt: pdfimages -list returned no images; the flowchart is vector page content.
% Grid evidence: tmp/review_2013_anhui_science/grid/page1_grid100.jpg and
% q2_grid50.jpg; no-grid source crop:
% tmp/review_2013_anhui_science/crops/q2_original_clean.png.
% Complete directed edges: start -> init(s=0,n=2) -> test(n<8?);
% test yes -> add(s=s+1/n) -> increment(n=n+2) -> test; test no -> output -> finish.
% The loop return joins the vertical test-entry line; no other edges exist.
\begin{tikzpicture}[
    baseline,
    >=Stealth,
    x=1cm,
    y=1cm,
    line width=0.45pt,
    line cap=round,
    line join=round,
    font=\normalsize,
    flowtext/.style={anchor=center,align=center,
        execute at begin node={\setlength{\parindent}{0pt}}},
    flowbox/.style={flowtext,draw,inner xsep=3pt,inner ysep=2pt},
]
  % Tight bounding box around the full source topology.
  \path[use as bounding box] (-2.95,-7.05) rectangle (4.35,0.42);

  \node[flowbox,rounded corners=0.16cm,minimum width=1.35cm,minimum height=0.68cm]
      (start) at (0,0) {开始};
  \node[flowbox,minimum width=2.35cm,minimum height=0.72cm]
      (init) at (0,-1.45) {$s=0,n=2$};
  \node[flowbox,diamond,aspect=3.0,minimum width=3.25cm,minimum height=1.04cm]
      (test) at (0,-3.20) {$n<8?$};
  \node[flowbox,minimum width=2.65cm,minimum height=1.22cm]
      (add) at (0,-4.95) {$s=s+\dfrac{1}{n}$};
  \node[flowbox,minimum width=2.55cm,minimum height=0.72cm]
      (inc) at (0,-6.50) {$n=n+2$};
  \node[flowbox,trapezium,trapezium left angle=75,trapezium right angle=105,
      minimum width=2.25cm,minimum height=0.82cm]
      (output) at (3.05,-4.95) {输出 $s$};
  \node[flowbox,rounded corners=0.16cm,minimum width=1.45cm,minimum height=0.72cm]
      (finish) at (3.05,-6.50) {结束};

  % Main path and the affirmative branch.
  \draw[->] (start.south) -- (init.north);
  \draw[->] (init.south) -- (test.north);
  \draw[->] (test.south) -- (add.north);
  \draw[->] (add.south) -- (inc.north);

  % The update box loops left and returns with one horizontal arrow into the
  % shared test-entry stem; there is no connection to the output branch.
  \coordinate (loopLeft) at (-2.62,-6.50);
  \coordinate (loopTop) at (-2.62,-2.20);
  \draw (inc.west) -- (loopLeft) -- (loopTop);
  \draw[->] (loopTop) -- (0,-2.20);

  % The negative branch runs independently around the right side to output.
  \coordinate (noRight) at (3.05,-3.20);
  \coordinate (noDrop) at (3.05,-4.50);
  \draw[->] (test.east) -- node[flowtext,above=2pt] {否} (noRight)
      -- (noDrop) -- (output.north);
  \draw[->] (output.south) -- (finish.north);

  \node[flowtext,inner sep=0pt] at (-0.55,-4.00) {是};
\end{tikzpicture}
